\documentclass[12pt,a4paper]{article}
\usepackage{fontspec}
\setmainfont{Times New Roman}[Path=C:/Windows/Fonts/,UprightFont=times.ttf,BoldFont=timesbd.ttf,ItalicFont=timesi.ttf,BoldItalicFont=timesbi.ttf]
\usepackage{amsmath,amssymb,amsfonts}
\usepackage[left=1.5cm,right=1.5cm,top=1.5cm,bottom=1.7cm,headheight=17pt]{geometry}
\usepackage{enumitem,array,fancyhdr,booktabs}

\pagestyle{fancy}
\fancyhf{}
\lhead{\small\textbf{ĐÁP ÁN ĐỀ CƯƠNG CHƯƠNG V -- PHƯƠNG PHÁP TỌA ĐỘ TRONG KHÔNG GIAN}}
\rhead{\small\textit{Biên soạn: Dương Gia Bảo}}
\rfoot{\small Trang \thepage}
\renewcommand{\headrulewidth}{0.4pt}
\fancypagestyle{firstpage}{\fancyhead{}\renewcommand{\headrulewidth}{0pt}}

\setlength{\parindent}{0pt}
\setlength{\parskip}{4pt}
\setlength{\emergencystretch}{1em}

\setlist[enumerate,1]{label=\textbf{Câu \arabic*:},leftmargin=*,itemsep=7pt,topsep=4pt}

\begin{document}
\thispagestyle{firstpage}

\noindent
\begin{minipage}[t]{0.36\textwidth}
    \centering
    {\fontsize{12}{14}\selectfont\textbf{ĐỀ CƯƠNG ÔN TẬP}}\\[4pt]
    \textit{Biên soạn: Dương Gia Bảo}\\[6pt]
    \setlength{\fboxrule}{0.8pt}
    \setlength{\fboxsep}{5pt}
    \fbox{\begin{minipage}{0.88\linewidth}
        \centering
        {\fontsize{12}{14}\selectfont\textbf{ĐÁP ÁN CHƯƠNG V}}
    \end{minipage}}
\end{minipage}\hfill
\begin{minipage}[t]{0.62\textwidth}
    \centering
    {\fontsize{13}{15}\selectfont\textbf{ĐÁP ÁN \& HƯỚNG DẪN GIẢI CHI TIẾT}}\\[3pt]
    {\fontsize{12}{14}\selectfont\textbf{CHƯƠNG V: PHƯƠNG PHÁP TỌA ĐỘ TRONG KHÔNG GIAN}}\\[3pt]
    {\fontsize{11}{13}\selectfont\textit{(GDPT 2018 -- 20 Trắc nghiệm, 5 Đúng/Sai, 25 Trả lời ngắn)}}
\end{minipage}

\par\vspace{4pt}
\hrule height 0.6pt
\vspace{8pt}

\begin{center}
    {\fontsize{13}{15}\selectfont\textbf{BẢNG ĐÁP ÁN NHANH}}
\end{center}

\subsubsection*{1. BẢNG ĐÁP ÁN PHẦN I -- TRẮC NGHIỆM NHIỀU LỰA CHỌN (20 CÂU)}

\begin{center}
\renewcommand{\arraystretch}{1.22}
\begin{tabular}{|c|c||c|c||c|c||c|c||c|c|}
\hline
\textbf{Câu} & \textbf{Đáp án} & \textbf{Câu} & \textbf{Đáp án} & \textbf{Câu} & \textbf{Đáp án} & \textbf{Câu} & \textbf{Đáp án} & \textbf{Câu} & \textbf{Đáp án} \\ \hline
\textbf{1} & \textbf{A} & \textbf{5} & \textbf{A} & \textbf{9} & \textbf{A} & \textbf{13} & \textbf{A} & \textbf{17} & \textbf{A} \\ \hline
\textbf{2} & \textbf{B} & \textbf{6} & \textbf{B} & \textbf{10} & \textbf{B} & \textbf{14} & \textbf{B} & \textbf{18} & \textbf{B} \\ \hline
\textbf{3} & \textbf{C} & \textbf{7} & \textbf{C} & \textbf{11} & \textbf{C} & \textbf{15} & \textbf{C} & \textbf{19} & \textbf{C} \\ \hline
\textbf{4} & \textbf{D} & \textbf{8} & \textbf{D} & \textbf{12} & \textbf{D} & \textbf{16} & \textbf{D} & \textbf{20} & \textbf{D} \\ \hline
\end{tabular}
\end{center}

\subsubsection*{2. BẢNG ĐÁP ÁN PHẦN II -- CÂU HỎI ĐÚNG SAI (5 CÂU)}

\begin{center}
\renewcommand{\arraystretch}{1.22}
\begin{tabular}{|c|c|c|c|c|}
\hline
\textbf{Câu hỏi} & \textbf{Lệnh a)} & \textbf{Lệnh b)} & \textbf{Lệnh c)} & \textbf{Lệnh d)} \\ \hline
\textbf{Câu 1} & \textbf{Đúng} & \textbf{Đúng} & \textbf{Đúng} & \textbf{Sai} \\ \hline
\textbf{Câu 2} & \textbf{Đúng} & \textbf{Sai} & \textbf{Đúng} & \textbf{Đúng} \\ \hline
\textbf{Câu 3} & \textbf{Đúng} & \textbf{Sai} & \textbf{Sai} & \textbf{Đúng} \\ \hline
\textbf{Câu 4} & \textbf{Đúng} & \textbf{Sai} & \textbf{Đúng} & \textbf{Đúng} \\ \hline
\textbf{Câu 5} & \textbf{Đúng} & \textbf{Đúng} & \textbf{Đúng} & \textbf{Sai} \\ \hline
\end{tabular}
\end{center}

\subsubsection*{3. BẢNG ĐÁP ÁN PHẦN III -- CÂU TRẢ LỜI NGẮN (25 CÂU)}

\begin{center}
\renewcommand{\arraystretch}{1.22}
\begin{tabular}{|c|c||c|c||c|c||c|c||c|c|}
\hline
\textbf{Câu} & \textbf{Đáp số} & \textbf{Câu} & \textbf{Đáp số} & \textbf{Câu} & \textbf{Đáp số} & \textbf{Câu} & \textbf{Đáp số} & \textbf{Câu} & \textbf{Đáp số} \\ \hline
\textbf{1} & $1$ & \textbf{6} & $2$ & \textbf{11} & $3$ & \textbf{16} & $1$ & \textbf{21} & $24$ \\ \hline
\textbf{2} & $3$ & \textbf{7} & $4$ & \textbf{12} & $4$ & \textbf{17} & $25$ & \textbf{22} & $-16$ \\ \hline
\textbf{3} & $2$ & \textbf{8} & $1$ & \textbf{13} & $90$ & \textbf{18} & $3$ & \textbf{23} & $3$ \\ \hline
\textbf{4} & $3$ & \textbf{9} & $2$ & \textbf{14} & $30$ & \textbf{19} & $5$ & \textbf{24} & $2$ \\ \hline
\textbf{5} & $2$ & \textbf{10} & $6$ & \textbf{15} & $17$ & \textbf{20} & $-11$ & \textbf{25} & $3$ \\ \hline
\end{tabular}
\end{center}

\vspace{8pt}
\hrule height 0.6pt
\vspace{8pt}

\begin{center}
    {\fontsize{13}{15}\selectfont\textbf{HƯỚNG DẪN GIẢI CHI TIẾT}}
\end{center}

\subsection*{PHẦN I. TRẮC NGHIỆM NHIỀU LỰA CHỌN (CÂU 1 -- 20)}

\begin{enumerate}
\item \textbf{Đáp án A.}  
Mặt phẳng $(P): 2x - 3y + z - 5 = 0$ có một vectơ pháp tuyến là $\vec{n}_1 = (2; -3; 1)$.

\item \textbf{Đáp án B.}  
Mặt phẳng đi qua $M(1; -2; 3)$ và nhận $\vec{n} = (3; 1; -2)$ làm VTPT có phương trình:
\[
3(x - 1) + 1(y + 2) - 2(z - 3) = 0 \iff 3x + y - 2z + 5 = 0.
\]

\item \textbf{Đáp án C.}  
Mặt phẳng đi qua ba điểm $A(2; 0; 0)$, $B(0; -3; 0)$, $C(0; 0; 4)$ có phương trình theo đoạn chắn là $\dfrac{x}{2} + \dfrac{y}{-3} + \dfrac{z}{4} = 1 \iff \dfrac{x}{2} - \dfrac{y}{3} + \dfrac{z}{4} = 1$.

\item \textbf{Đáp án D.}  
Khoảng cách từ điểm $M(1; 2; -1)$ đến mặt phẳng $(P): 2x - 2y + z + 12 = 0$:
\[
d(M, (P)) = \dfrac{|2(1) - 2(2) + 1(-1) + 12|}{\sqrt{2^2 + (-2)^2 + 1^2}} = \dfrac{|9|}{3} = 3.
\]

\item \textbf{Đáp án A.}  
Ta có $\dfrac{1}{2} = \dfrac{-2}{-4} = \dfrac{3}{6} \ne \dfrac{-1}{5}$. Do đó $(P) \parallel (Q)$.

\item \textbf{Đáp án B.}  
Đường thẳng $d: \dfrac{x-1}{2} = \dfrac{y+2}{-3} = \dfrac{z-5}{4}$ có một VTCP là $\vec{u}_2 = (2; -3; 4)$.

\item \textbf{Đáp án C.}  
Thay $t = 1$ vào phương trình tham số của $d$, ta được:
\[
\begin{cases} x = 1 + 2(1) = 3 \\ y = -1 - 1 = -2 \\ z = 3 + 4(1) = 7 \end{cases} \implies M(3; -2; 7) \in d.
\]

\item \textbf{Đáp án D.}  
Vì đường thẳng vuông góc với $(P): 3x - y + 2z - 1 = 0$ nên nhận $\vec{n} = (3; -1; 2)$ làm VTCP. Phương trình chính tắc đi qua $A(1; 0; -2)$ là $\dfrac{x-1}{3} = \dfrac{y}{-1} = \dfrac{z+2}{2}$.

\item \textbf{Đáp án A.}  
Đường thẳng $d$ có VTCP $\vec{u} = (1; 2; -1)$ và đi qua điểm $M(1; 2; 3)$.  
Mặt phẳng $(P)$ có VTPT $\vec{n} = (2; 1; 4)$.  
Ta có $\vec{u} \cdot \vec{n} = 1(2) + 2(1) + (-1)(4) = 0$.  
Thay tọa độ điểm $M(1; 2; 3)$ vào phương trình $(P)$: $2(1) + 2 + 4(3) - 16 = 0$ (thỏa mãn).  
Vậy đường thẳng $d$ nằm trong mặt phẳng $(P)$.

\item \textbf{Đáp án B.}  
Hình chiếu vuông góc của điểm $M(x_0; y_0; z_0)$ lên mặt phẳng tọa độ $(Oxy)$ có tọa độ là $(x_0; y_0; 0)$. Do đó $M'(2; -1; 0)$.

\item \textbf{Đáp án C.}  
Ta có $\vec{u}_1 \cdot \vec{u}_2 = 1(2) + 2(1) + 2(-2) = 0$.  
Vì tích vô hướng bằng $0$ nên góc giữa hai đường thẳng bằng $90^\circ$.

\item \textbf{Đáp án D.}  
VTPT của $(P)$ là $\vec{n}_1 = (1; 1; 0)$, VTPT của $(Q)$ là $\vec{n}_2 = (0; 1; 1)$.
\[
\cos((P), (Q)) = \dfrac{|\vec{n}_1 \cdot \vec{n}_2|}{|\vec{n}_1||\vec{n}_2|} = \dfrac{|1\cdot 0 + 1\cdot 1 + 0\cdot 1|}{\sqrt{1^2 + 1^2}\sqrt{1^2 + 1^2}} = \dfrac{1}{2} \implies ((P), (Q)) = 60^\circ.
\]

\item \textbf{Đáp án A.}  
VTCP của đường thẳng $d$ là $\vec{u} = (1; 1; \sqrt{2})$, VTPT của mặt phẳng $(Oxy)$ là $\vec{k} = (0; 0; 1)$.
\[
\sin(d, (Oxy)) = \dfrac{|\vec{u} \cdot \vec{k}|}{|\vec{u}||\vec{k}|} = \dfrac{|\sqrt{2}|}{\sqrt{1 + 1 + 2}\cdot 1} = \dfrac{\sqrt{2}}{2} \implies (d, (Oxy)) = 45^\circ.
\]

\item \textbf{Đáp án B.}  
\[
\cos(\vec{u}_1, \vec{u}_2) = \dfrac{|\vec{u}_1 \cdot \vec{u}_2|}{|\vec{u}_1||\vec{u}_2|} = \dfrac{|1(1) + (-1)(0) + 0(-1)|}{\sqrt{1 + 1}\sqrt{1 + 1}} = \dfrac{1}{2}.
\]

\item \textbf{Đáp án C.}  
Hai mặt phẳng vuông góc khi và chỉ khi tích vô hướng hai VTPT bằng $0$:
\[
\vec{n}_P \cdot \vec{n}_Q = 2(1) - 1(2) + 2m = 0 \iff 2m = 0 \iff m = 0.
\]

\item \textbf{Đáp án D.}  
Phương trình dạng $(x-a)^2 + (y-b)^2 + (z-c)^2 = R^2 \implies$ tâm $I(1; -2; 3)$ và bán kính $R = \sqrt{16} = 4$.

\item \textbf{Đáp án A.}  
Mặt cầu tâm $I(2; 1; -1)$ tiếp xúc $(Oxy): z = 0$ nên bán kính $R = d(I, (Oxy)) = |-1| = 1$.  
Phương trình mặt cầu là $(x-2)^2 + (y-1)^2 + (z+1)^2 = 1$.

\item \textbf{Đáp án B.}  
Ta có $a = 1, b = -2, c = 3, d = -2$.  
Bán kính $R = \sqrt{a^2 + b^2 + c^2 - d} = \sqrt{1^2 + (-2)^2 + 3^2 - (-2)} = \sqrt{16} = 4$.

\item \textbf{Đáp án C.}  
Mặt cầu $(S)$ có tâm $O(0; 0; 0)$ và bán kính $R = 3$.  
Khoảng cách từ $O$ đến $(P)$ là:
\[
d(O, (P)) = \dfrac{|-15|}{\sqrt{2^2 + (-1)^2 + 2^2}} = \dfrac{15}{3} = 5.
\]
Vì $d(O, (P)) = 5 > 3 = R$ nên mặt phẳng và mặt cầu không có điểm chung.

\item \textbf{Đáp án D.}  
Vùng không gian mặt cầu tâm $O(0; 0; 0)$ bán kính $R = 50\text{ km}$ có bề mặt giới hạn là $x^2 + y^2 + z^2 = 50^2 = 2500$.
\end{enumerate}

\subsection*{PHẦN II. TRẮC NGHIỆM ĐÚNG SAI (CÂU 1 -- 5)}

\begin{enumerate}
\item \textbf{Câu 1:}
\begin{enumerate}[label=\alph*)]
    \item \textbf{Đúng.} Điểm $A(2;0;0), B(0;3;0), C(0;0;6)$. Phương trình theo đoạn chắn là $\dfrac{x}{2} + \dfrac{y}{3} + \dfrac{z}{6} = 1 \iff 3x + 2y + z - 6 = 0$.
    \item \textbf{Đúng.} VTPT của $(ABC)$ là $\vec{n} = (3; 2; 1)$.
    \item \textbf{Đúng.} $d(O, (ABC)) = \dfrac{|-6|}{\sqrt{3^2 + 2^2 + 1^2}} = \dfrac{6}{\sqrt{14}} = \dfrac{3\sqrt{14}}{7}$.
    \item \textbf{Sai.} Thể tích khối tứ diện là $V = \dfrac{1}{6} OA \cdot OB \cdot OC = \dfrac{1}{6} \cdot 2 \cdot 3 \cdot 6 = 6$ (không phải $12$).
\end{enumerate}

\item \textbf{Câu 2:}
\begin{enumerate}[label=\alph*)]
    \item \textbf{Đúng.} Đường thẳng $d$ đi qua $M(2; -1; 0)$ và có VTCP $\vec{u} = (1; 2; -2)$.
    \item \textbf{Sai.} $(P)$ có VTPT $\vec{n} = (2; -1; 2)$. Vì $\dfrac{1}{2} \ne \dfrac{2}{-1}$ nên $\vec{u}$ không cùng phương với $\vec{n}$, do đó $d$ không vuông góc với $(P)$.
    \item \textbf{Đúng.} $(Q) \parallel (P)$ có dạng $2x - y + 2z + D = 0$. Vì $A(1; 2; -1) \in (Q)$ nên $2(1) - 2 + 2(-1) + D = 0 \iff D = 2$. Vậy $(Q): 2x - y + 2z + 2 = 0$.
    \item \textbf{Đúng.} Đường thẳng $\Delta$ qua $A(1; 2; -1)$ nhận $\vec{u} = (1; 2; -2)$ làm VTCP nên có phương trình chính tắc $\dfrac{x-1}{1} = \dfrac{y-2}{2} = \dfrac{z+1}{-2}$.
\end{enumerate}

\item \textbf{Câu 3:}
\begin{enumerate}[label=\alph*)]
    \item \textbf{Đúng.} $(S)$ có $a = 2, b = -1, c = 3, d = -11 \implies$ tâm $I(2; -1; 3)$ và bán kính $R = \sqrt{2^2 + (-1)^2 + 3^2 - (-11)} = \sqrt{25} = 5$.
    \item \textbf{Sai.} $d(I, (P)) = \dfrac{|2(2) + 2(-1) - 3 + 17|}{\sqrt{2^2 + 2^2 + (-1)^2}} = \dfrac{|16|}{3} = \dfrac{16}{3} \ne 4$.
    \item \textbf{Sai.} Vì $d(I, (P)) = \dfrac{16}{3} > 5 = R$ nên mặt phẳng $(P)$ không cắt mặt cầu $(S)$.
    \item \textbf{Đúng.} Mặt phẳng song song $(P)$ có dạng $2x + 2y - z + D = 0$. Tiếp xúc $(S) \iff d(I) = 5 \iff \dfrac{|2(2) + 2(-1) - 3 + D|}{3} = 5 \iff |D - 1| = 15 \iff D = 16$ hoặc $D = -14$.
\end{enumerate}

\item \textbf{Câu 4:}
\begin{enumerate}[label=\alph*)]
    \item \textbf{Đúng.} VTCP của $d_1$ là $\vec{u}_1 = (2; 1; -2)$, của $d_2$ là $\vec{u}_2 = (-1; 2; 2)$.
    \item \textbf{Sai.} $\vec{u}_1 \cdot \vec{u}_2 = 2(-1) + 1(2) + (-2)(2) = -4 \ne 0 \implies d_1$ không vuông góc với $d_2$.
    \item \textbf{Đúng.} $\cos(d_1, d_2) = \dfrac{|\vec{u}_1 \cdot \vec{u}_2|}{|\vec{u}_1||\vec{u}_2|} = \dfrac{|-4|}{3 \times 3} = \dfrac{4}{9}$.
    \item \textbf{Đúng.} Lấy $M_1(1; 0; -2) \in d_1$ và $M_2(-2; 1; 0) \in d_2 \implies \vec{M_1M_2} = (-3; 1; 2)$.  
    Ta có $[\vec{u}_1, \vec{u}_2] = (6; -2; 5) \implies [\vec{u}_1, \vec{u}_2] \cdot \vec{M_1M_2} = 6(-3) - 2(1) + 5(2) = -10 \ne 0$. Do đó $d_1$ và $d_2$ chéo nhau.
\end{enumerate}

\item \textbf{Câu 5:}
\begin{enumerate}[label=\alph*)]
    \item \textbf{Đúng.} Tọa độ $z_A = z_B = 3$, do đó cao độ của drone luôn bằng $3$ đơn vị $= 300\text{ m}$.
    \item \textbf{Đúng.} $\vec{AB} = (6; 8; 0) = 2(3; 4; 0)$, do đó $\vec{u} = (3; 4; 0)$ là một VTCP.
    \item \textbf{Đúng.} Độ dài $AB = \sqrt{6^2 + 8^2 + 0^2} = 10$ đơn vị $= 10 \times 100\text{ m} = 1000\text{ m}$.
    \item \textbf{Sai.} Hình chiếu của $AB$ lên mặt đất $(Oxy)$ là đoạn thẳng nối $(1; 2; 0)$ và $(7; 10; 0)$ nằm trên đường thẳng $4x - 3y + 2 = 0$. Khoảng cách từ $O(0;0;0)$ đến đường thẳng chứa $AB$ là $\sqrt{d(O, \Delta_{xy})^2 + 3^2} = \sqrt{\left(\dfrac{2}{5}\right)^2 + 9} = \sqrt{9{,}16} \approx 3{,}026$ đơn vị $\approx 302{,}6\text{ m} \ne 150\text{ m}$.
\end{enumerate}
\end{enumerate}

\subsection*{PHẦN III. TRẢ LỜI NGẮN (CÂU 1 -- 25)}

\begin{enumerate}
\item \textbf{Đáp số: $1$.}  
Trung điểm của $AB$ là $M(2; 1; 2)$. Vectơ $\vec{AB} = (2; -2; -2) = 2(1; -1; -1)$.  
Mặt phẳng trung trực nhận $\vec{n} = (1; -1; -1)$ làm VTPT:
\[
1(x - 2) - 1(y - 1) - 1(z - 2) = 0 \iff x - y - z + 1 = 0 \implies D = 1.
\]

\item \textbf{Đáp số: $3$.}  
$\vec{n}_Q = (1; -1; 1), \vec{n}_R = (2; 1; -1) \implies [\vec{n}_Q, \vec{n}_R] = (0; 3; 3)$.  
Mặt phẳng $(P)$ có VTPT $\vec{n}_P = (0; 3; 3)$ và đi qua $M(1; -2; 1)$:
\[
0(x - 1) + 3(y + 2) + 3(z - 1) = 0 \iff 3y + 3z + 3 = 0 \implies D = 3.
\]

\item \textbf{Đáp số: $2$.}  
\[
d(O, (P)) = \dfrac{|-14|}{\sqrt{2^2 + (-3)^2 + 6^2}} = \dfrac{14}{7} = 2.
\]

\item \textbf{Đáp số: $3$.}  
Lấy điểm $M(2; 0; 0) \in (P)$. Khoảng cách giữa hai mặt phẳng song song là:
\[
d((P), (Q)) = d(M, (Q)) = \dfrac{|2(2) - 0 + 2(0) + 5|}{\sqrt{2^2 + (-1)^2 + 2^2}} = \dfrac{9}{3} = 3.
\]

\item \textbf{Đáp số: $2$.}  
Điều kiện để hai mặt phẳng song song:
\[
\dfrac{m}{4} = \dfrac{2}{4} = \dfrac{-1}{-2} \ne \dfrac{1}{3} \iff \dfrac{m}{4} = \dfrac{1}{2} \iff m = 2.
\]

\item \textbf{Đáp số: $2$.}  
\[
d(S, (P)) = \dfrac{|2 + 2(3) + 2(4) - 10|}{\sqrt{1^2 + 2^2 + 2^2}} = \dfrac{|16 - 10|}{3} = \dfrac{6}{3} = 2\text{ mét}.
\]

\item \textbf{Đáp số: $4$.}  
Phương trình tham số của $d$: $x = 1 + 2t, y = -2 - t, z = 3 + t$.  
Thay vào $(P)$: $2(1 + 2t) + (-2 - t) - (3 + t) + 1 = 0 \iff 2t - 2 = 0 \iff t = 1$.  
Suy ra tọa độ giao điểm $M(3; -3; 4)$. Do đó $T = 3 + (-3) + 4 = 4$.

\item \textbf{Đáp số: $1$.}  
Đường thẳng đi qua $A(1; 3; -1)$ vuông góc với $(P)$ có phương trình: $x = 1 + t, y = 3 + 2t, z = -1 - t$.  
Thay vào $(P)$: $(1 + t) + 2(3 + 2t) - (-1 - t) - 2 = 0 \iff 6t + 6 = 0 \iff t = -1$.  
Tọa độ hình chiếu là $H(0; 1; 0) \implies S = 0 + 1 + 0 = 1$.

\item \textbf{Đáp số: $2$.}  
Hình chiếu của $A(1; -2; 3)$ lên $(Oyz)$ là $H(0; -2; 3) \implies A'(-1; -2; 3)$.  
Khoảng cách $AA' = \sqrt{(-1 - 1)^2 + 0^2 + 0^2} = 2$.

\item \textbf{Đáp số: $6$.}  
Xét hệ phương trình giao điểm của $d_1$ và $d_2$:
\[
\begin{cases} 2 + t = 2 + 2t' \\ 1 - t = 1 + t' \\ 3 + 2t = 3 - t' \end{cases} \iff \begin{cases} t = 0 \\ t' = 0 \end{cases}.
\]
Giao điểm là $I(2; 1; 3) \implies x_0 + y_0 + z_0 = 2 + 1 + 3 = 6$.

\item \textbf{Đáp số: $3$.}  
Đường thẳng $d$ đi qua $A(1; 2; 3)$ có VTCP $\vec{u} = (1; 2; 2)$. Ta có $\vec{AM} = (3; 3; 0)$.  
$[\vec{AM}, \vec{u}] = (6; -6; 3) = 3(2; -2; 1) \implies |[\vec{AM}, \vec{u}]| = 3 \times 3 = 9$.  
Khoảng cách là $d(M, d) = \dfrac{|[\vec{AM}, \vec{u}]|}{|\vec{u}|} = \dfrac{9}{3} = 3$.

\item \textbf{Đáp số: $4$.}  
$\vec{u}_1 = (1; 0; 0), \vec{u}_2 = (0; 1; 0) \implies [\vec{u}_1, \vec{u}_2] = (0; 0; 1)$.  
$\vec{AB} = (0; 2; -4)$.  
Khoảng cách ngắn nhất là $d(d_1, d_2) = \dfrac{|[\vec{u}_1, \vec{u}_2] \cdot \vec{AB}|}{|[\vec{u}_1, \vec{u}_2]|} = \dfrac{|-4|}{1} = 4\text{ km}$.

\item \textbf{Đáp số: $90$.}  
$\vec{u}_1 = (1; \sqrt{2}; 1)$ và $\vec{u}_2 = (1; -\sqrt{2}; 1)$.  
$\vec{u}_1 \cdot \vec{u}_2 = 1 - 2 + 1 = 0 \implies$ góc giữa hai đường thẳng bằng $90^\circ$.

\item \textbf{Đáp số: $30$.}  
$\vec{u} = (1; 0; 1)$, VTPT $\vec{n} = (1; 1; 0)$.  
$\sin\varphi = \dfrac{|1\cdot 1 + 0\cdot 1 + 1\cdot 0|}{\sqrt{2}\sqrt{2}} = \dfrac{1}{2} \implies \varphi = 30^\circ$.

\item \textbf{Đáp số: $17$.}  
$\vec{n}_P = (2; 2; -1) \implies |\vec{n}_P| = 3$; $\vec{n}_Q = (4; -3; 0) \implies |\vec{n}_Q| = 5$.  
$\cos\alpha = \dfrac{|2(4) + 2(-3) + (-1)(0)|}{3 \times 5} = \dfrac{2}{15} \implies a = 2, b = 15 \implies a + b = 17$.

\item \textbf{Đáp số: $1$.}  
$\vec{n}_P = (1; 0; 1) \implies |\vec{n}_P| = \sqrt{2}$; $\vec{n}_Q = (0; 1; m) \implies |\vec{n}_Q| = \sqrt{1 + m^2}$.  
$\cos 60^\circ = \dfrac{|m|}{\sqrt{2}\sqrt{1 + m^2}} = \dfrac{1}{2} \iff 2m^2 = 1 + m^2 \iff m^2 = 1 \implies m = 1$ (do $m > 0$).

\item \textbf{Đáp số: $25$.}  
$\vec{n}_\alpha = (3; 4; 12)$, $\vec{n}_{Oxy} = (0; 0; 1)$.  
$\cos\alpha = \dfrac{|12|}{\sqrt{3^2 + 4^2 + 12^2}\cdot 1} = \dfrac{12}{13} \implies a = 12, b = 13 \implies a + b = 25$.

\item \textbf{Đáp số: $3$.}  
Đoạn thẳng $AB = \sqrt{(3-1)^2 + (-2-2)^2 + (1 - (-3))^2} = \sqrt{4 + 16 + 16} = 6$.  
Bán kính mặt cầu là $R = \dfrac{AB}{2} = 3$.

\item \textbf{Đáp số: $5$.}  
Khoảng cách từ tâm $I(1; 2; 3)$ đến mặt phẳng $(P)$:
\[
d(I, (P)) = \dfrac{|2(1) - 2(2) + 3 + 8|}{\sqrt{2^2 + (-2)^2 + 1^2}} = \dfrac{9}{3} = 3.
\]
Bán kính mặt cầu là $R = \sqrt{d^2 + r^2} = \sqrt{3^2 + 4^2} = 5$.

\item \textbf{Đáp số: $-11$.}  
$a = 1, b = -2, c = 3$. Bán kính $R = \sqrt{a^2 + b^2 + c^2 - m} = \sqrt{1 + 4 + 9 - m} = \sqrt{14 - m}$.  
Theo bài ra $R = 5 \iff 14 - m = 25 \iff m = -11$.

\item \textbf{Đáp số: $24$.}  
Điểm $M(1 + 4t; 2 + 3t; 12) \in d$. Ta có $\vec{IM} = (4t; 3t; 9)$.  
Khoảng cách từ $I$ đến $d$: $IM^2 = (4t)^2 + (3t)^2 + 9^2 = 25t^2 + 81 \ge 81 \implies d(I, d) = 9\text{ km}$ (đạt tại $t = 0$).  
Độ dài quãng đường đi trong vùng phủ sóng:
\[
L = 2\sqrt{R^2 - d^2} = 2\sqrt{15^2 - 9^2} = 2\sqrt{144} = 24\text{ km}.
\]

\item \textbf{Đáp số: $-16$.}  
Tâm mặt cầu $I(1; 2; 3)$. Vectơ pháp tuyến của tiếp diện $(\alpha)$ là $\vec{n} = \vec{IM} = (1; 2; 2)$.  
Phương trình tiếp diện: $1(x - 2) + 2(y - 4) + 2(z - 5) = 0 \iff x + 2y + 2z - 20 = 0$.  
Suy ra $b = 2, c = 2, d = -20 \implies P = 2 + 2 - 20 = -16$.

\item \textbf{Đáp số: $3$.}  
Gọi $I$ là trung điểm của $AB \implies I(2; 3; 4)$.  
Ta có $MA^2 + MB^2 = 2MI^2 + \dfrac{AB^2}{2}$. Do đó $MA^2 + MB^2$ nhỏ nhất khi và chỉ khi $MI$ ngắn nhất, tức $M$ là hình chiếu vuông góc của $I$ lên $(P)$.  
Đường thẳng qua $I$ vuông góc với $(P)$ là $x = 2 + t, y = 3 + t, z = 4 + t$.  
Thay vào $(P)$: $(2 + t) + (3 + t) + (4 + t) - 3 = 0 \iff 3t + 6 = 0 \iff t = -2$.  
Tọa độ $M(0; 1; 2) \implies x_M + y_M + z_M = 0 + 1 + 2 = 3$.

\item \textbf{Đáp số: $2$.}  
Gọi $H$ là hình chiếu vuông góc của $A$ lên đường thẳng $d$. Ta luôn có $d(A, (P)) \le AH$, dấu đẳng thức xảy ra khi $(P) \perp AH$.  
Do đó khoảng cách lớn nhất bằng khoảng cách từ $A$ đến đường thẳng $d$.  
Đường thẳng $d$ đi qua $B(1; -1; 2)$ có VTCP $\vec{u} = (1; 2; 2)$. Ta có $\vec{BA} = (2; 2; 0)$.  
$[\vec{BA}, \vec{u}] = (4; -4; 2) = 2(2; -2; 1) \implies |[\vec{BA}, \vec{u}]| = 2 \times 3 = 6$.  
Khoảng cách lớn nhất là $d_{\max} = \dfrac{|[\vec{BA}, \vec{u}]|}{|\vec{u}|} = \dfrac{6}{3} = 2$.

\item \textbf{Đáp số: $3$.}  
Khoảng cách giữa hai drone là độ dài đoạn $D_1D_2$:
\[
\vec{D_1D_2} = (t + 3 - 2t; 2t - 3 - t; 2 - 0) = (3 - t; t - 3; 2).
\]
$D_1D_2^2 = (3 - t)^2 + (t - 3)^2 + 2^2 = 2(t - 3)^2 + 4 \ge 4$.  
Dấu đẳng thức xảy ra khi $t - 3 = 0 \iff t = 3$ phút.
\end{enumerate}

\end{document}
